% ============================================================================
%  Subdocumento 13. Contenido. Se usa igual en el documento único y en el
%  PDF propio del subdocumento. Los formularios que le asigna el T-21 se
%  agregan solos al final (datos en formularios/ de esta carpeta).
% ============================================================================

\begin{resumenApertura}
\marcador{Qué resuelve este subdocumento, en cuatro o cinco líneas}
\begin{recibe}
  \item \marcador{Qué recibe el mandante}
\end{recibe}
\end{resumenApertura}

\section{Cartera de innovaciones}\label{sec:13-cartera-de-innovaciones}

\marcador{Las cinco innovaciones y su pertinencia al Caso}

\section{Tipo 1. Producto o servicio}\label{sec:13-tipo-1-producto-o-servicio}

\marcador{Resumen y análisis. La ficha va en el Formulario T-19}

\section{Tipo 2. Proceso}\label{sec:13-tipo-2-proceso}

\marcador{Resumen y análisis. La ficha va en el Formulario T-19}

\section{Tipo 3. Tecnológica o de arquitectura}\label{sec:13-tipo-3-tecnologica-o-de-arquitec}

\marcador{Resumen y análisis. La ficha va en el Formulario T-19}

\section{Tipo 4. Modelo de negocio o de contratación}\label{sec:13-tipo-4-modelo-de-negocio-o-de-co}

\marcador{Resumen y análisis. La ficha va en el Formulario T-19}

\section{Tipo 5. Experiencia de usuario, sostenibilidad o impacto social}\label{sec:13-tipo-5-experiencia-de-usuario-so}

\marcador{Resumen y análisis. La ficha va en el Formulario T-19}

\section{Trazabilidad con arquitectura, EDT y flujo de caja}\label{sec:13-trazabilidad-con-arquitectura-ed}

\marcador{Dónde se inserta cada innovación}
